\begin{enumerate}
    \item \textbf{Bernoulli 数(A027641)}

          见 \fullref{ssec:sheet-sequences-bernoulli}.

    \item \textbf{四个柱子的汉诺塔(A007664)}

          0, 1, 3, 5, 9, 13, 17, 25, 33, 41, 49, 65, 81, 97, 113, 129, 161, 193, 225, 257, 289, 321, 385, 449, \dots

          差分之后可以发现其实就是 1 次 +1, 2 次 +2, 3 次 +4, 4 次 +8\dots 的规律.

    \item \textbf{Ulam 数(A002858)}

          1, 2, 3, 4, 6, 8, 11, 13, 16, 18, 26, 28, 36, 38, 47, 48, 53, 57, 62, 69, 72, 77, 82, 87, 97, 99, 102, 106, 114, 126, 131, 138, 145, 148, 155, 175, 177, 180, 182, 189, 197, 206, 209, 219, 221, 236, 238, 241, 243, 253, 258, 260, 273, 282, 309, 316, 319, 324, 339, \dots

          \( a_1 = 1,\; a_2 = 2 \), \(a_n\) 表示在所有 \(>a_{n-1}\) 的数中, 最小的, 能被表示成 (前面的两个不同的元素的和) 的数.
\end{enumerate}
